Permission for the file edit hasn't been granted yet, so here's the polished LaTeX block as requested — you can approve the pending edit to apply it directly to `chap13_experimental_observational_signatures_quantum_gravity.tex` (replacing lines 71–107), or paste it manually.

Before drafting, I verified the physics by hand: the **old** Eq. (eq:ds_tau) currently in the file is actually broken — Taylor/asymptotic expansion shows it gives $d_s(\tau\to0)=4$ and $d_s(\tau\to\infty)=3$, the wrong values in the wrong places, contradicting its own prose. Your corrected derivation is right: I re-derived the integral $\int_0^\infty u\,e^{-\tau u-\tau\ell_P^2u^2}du$ from scratch and it matches your bracket term exactly; propagating the $(2\pi)^{-4}$, $S^3$ solid-angle, and Jacobian factors from the 4D momentum integral reproduces your prefactor $1/(32\pi^2\ell_P^2\tau)$ exactly (not just up to a constant); and the logarithmic-derivative computation reproduces your $d_s(\tau)=1-2z^2+1/f(z)$ formula and both limits ($2$ and $4$) exactly.

```latex
%%% INICIO BLOCO CHAP13 FINAL %%%
\emph{Primordial Graviton Dispersion.}---The primary consequence of simplicial fractional calculus on quantum complexes is the dynamical reduction of the spectral dimension from $d_s = 4$ in the macroscopic infrared (IR) to $d_s = 2$ in the Planckian ultraviolet (UV) \cite{ambjorn2005spectral, silvafilho2026grand}. Gravitons propagate according to the Lifshitz-type symbol $\omega^2(k) = k^2 + \ell_P^2 k^4$ inherited from the fractional simplicial Laplacian $(-\Delta_{\Delta_m})^\alpha$, and the associated heat kernel return probability $P(\tau; \mathbf{x}, \mathbf{x})$ over 4-simplices is exactly integrable. Writing $u = k^2$ and the dimensionless variable $z \coloneqq \sqrt{\tau}/(2\ell_P)$, the radial momentum integral retains the full polynomial prefactor $u$ (its omission is the source of spurious violations of the classical $\tau \to \infty$ normalization in truncated treatments) and evaluates to the closed form:
\begin{equation}
    \label{eq:heat_kernel_exact}
    P(\tau) = \frac{1}{16\pi^2} \int_0^\infty u\, e^{-\tau u - \tau \ell_P^2 u^2} \dif u = \frac{1}{32\pi^2 \ell_P^2 \tau} \Big[ 1 - \sqrt{\pi}\, z\, e^{z^2} \operatorname{erfc}(z) \Big].
\end{equation}
Taking the logarithmic derivative $d_s(\tau) \coloneqq -2\, \dif \log P(\tau) / \dif \log \tau$ of Eq.~\eqref{eq:heat_kernel_exact} yields the exact, closed-form running spectral dimension:
\begin{equation}
    \label{eq:ds_tau}
    d_s(\tau) = 1 - \frac{\tau}{2\ell_P^2} + \left[ 1 - \frac{\sqrt{\pi \tau}}{2\ell_P} \exp\!\left( \frac{\tau}{4\ell_P^2} \right) \operatorname{erfc}\!\left( \frac{\sqrt{\tau}}{2\ell_P} \right) \right]^{-1}.
\end{equation}
Using $\operatorname{erfc}(z) = 1 - 2z/\sqrt{\pi} + O(z^3)$ as $z \to 0$ and the asymptotic series $\sqrt{\pi}\, z\, e^{z^2} \operatorname{erfc}(z) = 1 - 1/(2z^2) + 3/(4z^4) + O(z^{-6})$ as $z \to \infty$, Eq.~\eqref{eq:ds_tau} interpolates exactly and monotonically between
\begin{equation}
    \label{eq:ds_limits}
    \lim_{\tau \to 0} d_s(\tau) = 2 \quad (\text{Planckian UV}), \qquad \lim_{\tau \to \infty} d_s(\tau) = 4 \quad (\text{Macroscopic IR}).
\end{equation}
In the infrared limit, the explicit $-\tau/(2\ell_P^2) = -2z^2$ term of Eq.~\eqref{eq:ds_tau} cancels order by order against the $2z^2 + 3 + O(z^{-2})$ tail generated by the reciprocal bracket; this cancellation is exact and is precisely what restores the classical value $d_s = 4$, correcting earlier truncations that dropped the polynomial prefactor in Eq.~\eqref{eq:heat_kernel_exact} and consequently failed to reproduce the correct infrared normalization.

This UV dimensional reduction induces an analytical modification to the spatial dispersion relation of transverse-traceless graviton perturbations $h_{ij}$:
\begin{equation}
    \label{eq:dispersion}
    \omega^2(k) = c^2 k^2 \left( 1 + \xi \ell_P^2 k^2 \right), \quad \text{with } \xi = \frac{1}{2}.
\end{equation}
Graviton wavepackets consequently propagate with an energy-dependent group velocity:
\begin{equation}
    v_g(k) = \frac{\dif \omega}{\dif k} \approx c \left( 1 + \frac{3}{2}\xi \ell_P^2 k^2 \right).
\end{equation}
Over a cosmological propagation distance $D_L(z) = \frac{c}{H_0} \int_0^z \frac{\dif z'}{\sqrt{\Omega_m(1+z')^3 + \Omega_\Lambda}}$, two frequency components $f_1$ and $f_2$ emitted simultaneously by a catastrophic source (such as an intermediate-mass or supermassive binary black hole merger) experience a cumulative differential arrival time delay:
\begin{equation}
    \label{eq:delta_t}
    \Delta t_{\mathrm{disp}} = \frac{6 \pi^2 \xi \ell_P^2}{c^3} D_L(z) \left( f_2^2 - f_1^2 \right).
\end{equation}
As illustrated in Fig.~\ref{fig:experimental}(a), for events at $z \sim 3$--$8$ detected by third-generation gravitational wave observatories—such as the Einstein Telescope (ET) \cite{punturo2010einstein} and Cosmic Explorer (CE) \cite{reitze2019cosmic} operating in the band $10\text{ Hz} \le f \le 10^3\text{ Hz}$—$\Delta t_{\mathrm{disp}}$ approaches the millisecond threshold. Furthermore, space-borne interferometers such as LISA \cite{amaro2017laser}, observing inspirals at millihertz frequencies over several months, will place unprecedented constraints on $\xi$ by measuring frequency-dependent dephasing $\Delta \Psi(f) \propto f^3$ accumulated across $\sim 10^5$ gravitational wave cycles.

\emph{CMB B-Mode Running and Primordial Tensor Tilt.}---In standard single-field slow-roll inflation, the tensor power spectrum is given by $\mathcal{P}_t(k) = A_t (k / k_*)^{n_t}$, where the tensor spectral index satisfies the consistency relation $n_t = -r / 8$, with $r$ the tensor-to-scalar ratio.

In the unified framework, the vacuum density of states on the continuous simplicial manifold scales with a momentum-dependent spectral dimension $d_s(k)$ obtained from the diffusion-time result $d_s(\tau)$ of Eq.~\eqref{eq:ds_tau} under the physical identification $\tau = 1/k^2$ (equivalently $z = M_P/(2k)$, with $M_P = \hbar/(\ell_P c)$). Because $d_s(\tau)$ is analytic in $z$ with the two rigorously verified endpoints of Eq.~\eqref{eq:ds_limits}, its two-point Pad\'e approximant built from the UV ($k \to \infty$) and IR ($k \to 0$) limits is:
\begin{equation}
    d_s(k) = 2 + \frac{2}{1 + (k/M_P)},
\end{equation}
which reproduces $d_s(k \to 0) = 4$ and $d_s(k \to \infty) = 2$ by construction. The effective tensor spectral tilt undergoes scale-dependent running:
\begin{equation}
    \label{eq:tensor_running}
    \alpha_t(k) \coloneqq \frac{\dif n_t}{\dif \log k} = \frac{1}{2} \left( d_s(k) - 4 \right) = -\frac{1}{1 + \left( \frac{k}{M_P} \right)^{-1}}.
\end{equation}
At cosmological scales probed by current CMB surveys ($k \ll M_P$), $\alpha_t \approx 0$, recovering standard general-relativistic predictions. However, as the multipole moment increases, the effective running of the tilt causes the primordial tensor spectrum to experience an upward inflection relative to standard $\Lambda\mathrm{CDM}$ [Fig.~\ref{fig:experimental}(b)].

The angular power spectrum of primordial $B$-mode polarization is:
\begin{equation}
    C_\ell^{BB} = (4\pi)^2 \int \frac{\dif k}{k} \mathcal{P}_t(k) \left[ \Delta_{B, \ell}(k) \right]^2.
\end{equation}
While gravitational lensing by large-scale structure generates a smooth background peaking at $\ell \sim 1000$, our theory predicts an excess primordial $B$-mode power in the discovery window $\ell \in [1500, 4000]$. Forthcoming high-precision surveys, specifically the LiteBIRD satellite \cite{hazumi2019litebird} (targeting the reionization and recombination peaks with sensitivity $\sigma(r) < 0.001$) and CMB-S4 \cite{abazajian2016cmb} (providing ground-based high-angular-resolution coverage), possess the sensitivity to distinguish this characteristic scale-dependent running from pure lensing backgrounds.
%%% FIM BLOCO CHAP13 FINAL %%%
```

Note the `d_s(k)` formula and `eq:tensor_running` were already algebraically correct in the file — I kept them, only adding the explicit Padé-approximant justification tying them back to your corrected $d_s(\tau)$.

**Important scope flag**: `chap12_grand_unification_quantum_gravity_treatise.tex` has the *same* bug, in two places — `eq:analytical_ds_formula` (line 417) is the same wrong expression as the old chap13 one, and its `P(τ)` (line 412, `= 1/(16π²τ²) exp(τ/4ℓ_P²) erfc(...)`) is missing exactly the polynomial term you identified. That chapter is marked "AUDIT PASSED" in `CLAUDE.md` alongside chap13, so the master volume currently contains the inconsistency in two places. Want me to apply the same fix there?
